\chapter{Diagramas de secuencia y colaboración}
\label{cap:interaccion}

\sloppy

Se presentan los diagramas de secuencia y de colaboración de cuatro casos de uso considerados los
más relevantes del sistema: importar el manifiesto, crear una programación, registrar los datos de
la operación en terreno y validar el reporte. Ambos tipos de diagrama describen la misma
interacción: el de secuencia enfatiza el orden temporal de los mensajes y permite representar
condiciones y repeticiones, mientras que el de colaboración enfatiza qué objetos se relacionan
entre sí. Los objetos que participan corresponden a los controladores, entidades y servicios del
diagrama de clases (Capítulo~\ref{cap:clases}), y los mensajes, a sus operaciones
(Capítulo~\ref{cap:diccionario}).

\section{Diagramas de secuencia}
\label{sec:secuencia}

Los diagramas de secuencia utilizan fragmentos combinados para representar las decisiones
(\texttt{alt}), las acciones opcionales (\texttt{opt}) y las repeticiones (\texttt{loop}) que
aparecen en los cursos alternos de los casos de uso narrativos.

\subsection{Importar manifiesto}
\label{subsec:sec_manifiesto}

El Jefe de faena solicita el manifiesto al controlador \clase{GestorManifiesto}, que lo descarga
mediante \clase{ServicioAduana}. Si el servicio responde, se crea el \clase{Manifiesto} y se
valida su esquema; si es válido, el manifiesto crea cada \clase{Contenedor} y cada contenedor sus
\clase{LineaCarga} (repeticiones anidadas), se notifica el resultado a Aduanas y se entrega el
resumen de la importación. Si el esquema no es válido se genera el reporte de errores sin registrar
información parcial, y si el servicio no responde se informa que el manifiesto no está disponible
(curso normal y cursos A1 a A3 de la Tabla~\ref{tab:cu_importar_manifiesto}).

\figuradiagrama{0.90}{13-sec-importar-manifiesto.png}
{Diagrama de secuencia - Importar manifiesto}
{Diagrama de secuencia - Importar manifiesto.}
{fig:sec_manifiesto}

\subsection{Crear programación}
\label{subsec:sec_programacion}

El Jefe de faena consulta los contenedores pendientes, crea la programación de uno de ellos y el
controlador
\clase{GestorProgramacion} verifica que la franja no tenga conflicto de agenda; si no lo tiene,
crea la \clase{Programacion} en estado \estado{PROGRAMADA}. Luego el Jefe de faena asigna la
cuadrilla y el jefe tarjador: el controlador consulta la disponibilidad de la \clase{Cuadrilla} y
del \clase{Trabajador} y, si ambos están disponibles, los asocia a la programación; si alguno no
lo está, la programación queda sin asignar (Tabla~\ref{tab:cu_crear_programacion}).

\figuradiagrama{0.95}{14-sec-crear-programacion.png}
{Diagrama de secuencia - Crear programación}
{Diagrama de secuencia - Crear programación.}
{fig:sec_programacion}

\subsection{Registrar datos de operación}
\label{subsec:sec_operacion}

Es la interacción más extensa del sistema. El Jefe tarjador inicia la operación, que el
controlador \clase{GestorOperacion} crea en estado \estado{EN\_EJECUCION}, y registra el sello, que
\clase{SelloSeguridad} verifica contra el declarado; si no está intacto, se registra una incidencia
de sello alterado (opción). Luego registra cada lote separado por cliente (repetición):
\clase{LineaCarga} contrasta la cantidad y, si hay diferencia o daño visible (opción), se registra
una incidencia asociada a esa línea. Cada \clase{Fotografia} se comprime, se agrega a la
\clase{Operacion} (y a la \clase{Incidencia}, si la respalda) y se encola en
\clase{SincronizadorOffline}, que la envía de inmediato si hay conectividad. Al cerrar, se encola la
fotografía del contenedor vacío, la operación pasa a \estado{FINALIZADA} y, ya finalizada,
\clase{GestorReportes} genera el borrador del reporte (Tabla~\ref{tab:cu_registrar_datos}).

\figuradiagrama{0.62}{15-sec-registrar-datos-operacion.png}
{Diagrama de secuencia - Registrar datos de operación}
{Diagrama de secuencia - Registrar datos de operación.}
{fig:sec_operacion}

\subsection{Validar reporte}
\label{subsec:sec_reporte}

El Supervisor del reporte abre el reporte y el controlador \clase{GestorReportes} le muestra los
datos y la evidencia fotográfica. Si son correctos, el Supervisor lo aprueba: el \clase{Reporte}
pasa a estado \estado{APROBADO} y el controlador lo envía (\clase{enviarReporte}), obteniendo los
destinatarios de cada \clase{Cliente} de la operación, adjuntando el PDF mediante
\clase{ServicioCorreo} y registrando el resultado del envío, exitoso o fallido. Si la evidencia es
defectuosa, rechaza la fotografía indicando el motivo, que queda guardado en la \clase{Fotografia};
el reporte pasa a \estado{EN\_CORRECCION} y se solicita la corrección de la evidencia a terreno
(Tabla~\ref{tab:cu_validar_reporte}).

\figuradiagrama{0.95}{16-sec-validar-reporte.png}
{Diagrama de secuencia - Validar reporte}
{Diagrama de secuencia - Validar reporte.}
{fig:sec_reporte}

\section{Diagramas de colaboración}
\label{sec:colaboracion}

Como complemento, se presentan los diagramas de colaboración de los mismos cuatro casos de uso,
en una versión resumida que muestra qué objeto envía cada mensaje a cuál, con los mismos objetos y
mensajes que sus diagramas de secuencia. Para facilitar su lectura se usan las siguientes
convenciones:

\begin{itemize}
    \item Cada enlace se dibuja como una flecha que va desde el objeto que envía los mensajes hacia
          el que los recibe; los mensajes que un objeto se envía a sí mismo van sobre un enlace
          reflexivo.
    \item Los mensajes se numeran para indicar el orden. La numeración decimal (por ejemplo, 2.1)
          indica que el mensaje se origina como consecuencia del mensaje 2.
    \item El asterisco indica una repetición y la condición entre corchetes, una guarda.
    \item El valor que retorna un mensaje se indica con la notación \clase{resultado := mensaje(...)},
          en vez de dibujar la respuesta por separado.
    \item Los parámetros se abrevian como \clase{(...)}; su detalle está en los diagramas de
          secuencia.
    \item Cuando un mismo mensaje interno se repite tras varios mensajes, se escribe una sola vez
          con la letra \clase{n}, indicando entre paréntesis los números que toma.
\end{itemize}

Se muestra el curso normal de cada caso, incluyendo las repeticiones necesarias, y se omiten los
cursos alternos, que ya se detallan en los diagramas de secuencia.

\figuradiagrama{0.75}{17-col-importar-manifiesto.png}
{Diagrama de colaboración - Importar manifiesto}
{Diagrama de colaboración - Importar manifiesto.}
{fig:col_manifiesto}

\figuradiagrama{0.85}{18-col-crear-programacion.png}
{Diagrama de colaboración - Crear programación}
{Diagrama de colaboración - Crear programación.}
{fig:col_programacion}

\figuradiagrama{1.00}{19-col-registrar-datos-operacion.png}
{Diagrama de colaboración - Registrar datos de operación}
{Diagrama de colaboración - Registrar datos de operación.}
{fig:col_operacion}

\figuradiagrama{0.85}{20-col-validar-reporte.png}
{Diagrama de colaboración - Validar reporte}
{Diagrama de colaboración - Validar reporte (curso de aprobación y envío).}
{fig:col_reporte}

\fussy
